\section{Motivation and Problem Statement}
\label{sec:intro-motivation}

Virtual Knowledge Graphs (VKGs) let a relational database be queried as if it were RDF, through an
Ontology-Based Data Access (OBDA) mapping layer such as Ontop~\cite{ontop}. Combined with an
LLM-driven natural-language-to-SPARQL (NL2SPARQL) pipeline, this lets a non-expert user ask a
question in plain language --- \emph{``Churches in Verona''} --- and receive results retrieved
directly from the underlying data, with no SPARQL expertise required. This is the architecture
underlying the Verona tourism domain used as one of this thesis's evaluation
domains~\cite{balasooriya2026}.

Prior work describing this architecture explicitly names \emph{``recommendations''} as future
work~\cite{balasooriya2026}: the system answers a question but does nothing to help a user
formulate a better next one. This thesis addresses that gap along three fronts, in increasing
order of research weight:

\begin{enumerate}[label=(\arabic*)]
    \item \textbf{Domain selection} --- generalize the application so a user can choose \emph{which}
        VKG to query (e.g.\ Movie vs.\ Tourism) instead of the system being wired to a single
        dataset.
    \item \textbf{User selection instead of login} --- let the user pick \emph{who they are} from a
        list of personas, so the system can reason about that user's preferences, without
        implementing real authentication.
    \item \textbf{Query recommendation} --- the thesis's core contribution: after answering a
        query, propose three follow-up query suggestions. If a user asks \emph{``Churches in
        Verona,''} the system should show results \emph{plus} suggestions such as \emph{``Roman
        Churches in Verona,''} \emph{``Free Churches in Verona,''} or \emph{``Trip to visit churches
        in Verona.''} Critically, this recommender should consider not only the user's previous
        queries but also their \textbf{real interaction data} --- because what a user \emph{says}
        they want and what they \emph{actually} engage with are frequently different.
\end{enumerate}

Item (3) is this thesis's actual contribution, and its framing is the thesis's central research
question: can a recommender that conditions on both a user's stated queries \emph{and} their real
interaction behavior produce suggestions that are measurably better --- more aligned with what the
user actually wants --- than one conditioning on query history alone?

\section{Scope of This Thesis}
\label{sec:intro-scope}

After items (1) and (2) had already been partially explored, the scope was narrowed to item (3).
Domain selection and user selection instead of login are handled by a separate, existing
application that has no query-recommendation functionality of its own. This thesis does not build
or integrate with that application's user interface. Its entire contribution is item (3): the
query recommender, delivered as a \textbf{standalone model and API} --- any calling application
integrates it later by calling an endpoint. This thesis is not responsible for that integration.

The multi-domain VKG infrastructure that had already been built at the time of this scope decision
--- two working VKGs (Movie and Tourism) and a domain-switching development harness --- is retained,
but strictly as an \textbf{internal development and evaluation harness}: a way to build, test, and
demo the recommender against real schemas and data, without depending on the external
application's timeline or codebase. It is documented in Chapter~\ref{ch:arch} because the
recommender's schema-grounding step and one of its two evaluation domains depend on it, but it is
explicitly not, itself, a thesis deliverable.

\section{Contributions}
\label{sec:intro-contrib}

This thesis makes the following contributions:

\begin{itemize}
    \item A \textbf{per-user behavioral preference profile} (Chapter~\ref{ch:behavior}) that
        quantifies, for the Movie domain, how far a user's stated search intent diverges from what
        they actually watch --- including an engineered temporal-join heuristic recovering an
        implicit link between free-text search queries and later watch events that does not exist
        as a key in the source data.
    \item A \textbf{retrieval-augmented, behavior-conditioned query recommender}
        (Chapter~\ref{ch:recommender}) --- a four-stage pipeline (retrieval, LLM-based candidate
        generation, schema-grounding validation, and behavior-aware ranking) delivered as a
        standalone HTTP API, adapting the GQR/RA-GQR line of work~\cite{bacciu2024} and borrowing
        scoring-axis ideas from classical, non-LLM query-recommendation research~\cite{barman2024}.
    \item An \textbf{empirically-grounded negative result} on schema-grounding validation
        (Section~\ref{sec:rec-validation}): a precision hallucination filter, motivated directly by
        a documented risk in the LLM+KG integration literature~\cite{pan2024} and by the baseline
        system's own measured SPARQL-correctness rate~\cite{balasooriya2026}, was designed,
        calibrated, and found not to be achievable with acceptable precision given how sparse this
        project's source ontologies are --- a concrete, reportable limitation rather than a silently
        shipped, over-claimed feature.
    \item A \textbf{controlled ablation study} (Chapter~\ref{ch:eval}) isolating the effect of the
        behavioral signal on suggestion quality, scored against real held-out ground truth, showing
        a $6.5\times$ increase in behavioral-alignment@3 and a more than $2\times$ increase in
        proxy-recall@3 when the behavioral profile is included --- the thesis's central empirical
        confirmation that conditioning on real behavior, not just stated query text, produces
        measurably different and better-aligned recommendations.
\end{itemize}

\section{Thesis Structure}
\label{sec:intro-outline}

Chapter~\ref{ch:background} introduces the theoretical background required for this work, including
Knowledge Graphs, Virtual Knowledge Graphs, Ontology-Based Data Access, Ontop, RDF, SPARQL, and
Large Language Models. Chapter~\ref{ch:relwork} reviews the four bodies of literature that ground
the recommender's design. Chapter~\ref{ch:arch} describes the system architecture and development
harness the recommender is built and evaluated against. Chapter~\ref{ch:behavior} presents the
behavioral preference profile, the empirical substrate for the thesis's central premise.
Chapter~\ref{ch:recommender} presents the recommender core itself --- the thesis's central design
contribution. Chapter~\ref{ch:eval} presents the evaluation and its results.
Chapter~\ref{ch:conclusion} summarizes the contributions, discusses limitations, and outlines
future work.
